%% kapa: exam=1 round=28 year=2017 period=1 form=A subject=civil_law subject_ko=민법 seq=12 number=12
%% source: Q-Net 감정평가사 기출문제 (cst00309) artlSeq=5199603 file=extracted/1008/5199603/2017년 제28회 감정평가사 1차시험 문제지/1교시 A형.pdf
%% conversion: kordoc 4.24.51 -> tools/gen1.py (auto); review=proofread-sonnet; flags=none
\begin{question}{12}
민법상 임의대리에 관한 설명으로 옳지 않은 것은?
\begin{choices}
  \item 대리인은 행위능력자임을 요하지 않는다.
  \item 대리권은 다른 특약이 없으면 법률관계의 종료 전에 수권행위를 철회한 경우에도 소멸한다.
  \item 대리인이 그 권한 내에서 본인을 위한 것임을 표시한 의사표시는 직접 본인에 대하여 효력이 생긴다.
  \item 특정한 법률행위를 위임한 경우에 대리인이 본인의 지시에 좇아 그 행위를 한 때에는 본인은 자기가 안 사정 또는 과실로 인하여 알지 못한 사정에 관하여 대리인의 부지를 주장하지 못한다.
  \item 대리인이 수인(數人)인 경우에 대리인은 원칙적으로 공동으로 대리하고 수권행위 또는 법률로 달리 정하는 경우에만 각자 본인을 대리한다.
\end{choices}
\end{question}
\answer{5}
